\begin{abstract}
	Low-light image enhancement is not only a brightness correction problem: an effective model must recover structure, preserve colors, control artifacts, and remain reliable when the test distribution differs from the training one. This work studies this problem under a cross-dataset generalization protocol. All models are trained on LOL-v2 Real and evaluated on the held-out LOL-v2 Real test set and on three independent external datasets: LOL-v2 Synthetic, MILLs, and ExDark. LOL-v1 is additionally reported as a non-independent overlap-control set: it was retained for descriptive comparison because its low-light inputs overlap with LOL-v2 Real, but it is excluded from claims of independent cross-dataset generalization. The evaluation combines full-reference metrics, no-reference metrics, diagnostic indicators, and qualitative failure analysis. The experimental workflow includes deterministic preprocessing, validation-based model selection, multi-seed evaluation, and a total of 624 logged experimental runs. Four model families are compared: (i) a classical U-Net, (ii) an illumination-guided U-Net, (iii) Frequency-Aware Exposure-Consistent Low-Light Enhancement (FAC-ECL), and (iv) Exposure-Consistency-based Failure-Aware Refinement (ECL-FAR). The results show that simply providing an explicit illumination cue can reduce residual darkness, but does not guarantee better reconstruction or color stability. FAC-ECL improves paired reconstruction most clearly on MILLs through frequency-aware supervision and exposure consistency. ECL-FAR, an original refinement-based model family introduced in this project, uses photometric instability as a cue for local correction. Although it is only a first implementation and was not exhaustively optimized in all its components, ECL-FAR is already competitive with FAC-ECL, achieves the highest observed SSIM on MILLs, and is visually convincing in several of the final qualitative comparisons.
\end{abstract}